\begin{claim}{C.32}{Proved}{Discrete proper operator groups on open orbits}
Let $O$ be a connected noncompact second-countable surface without
boundary, and let a discrete group $Q\leq\operatorname{Homeo}(O)$ act
freely and properly. With compact-open topology on the homeomorphism
groups, inclusion induces an injection
\[
 \pi_1\!\left(C_{\operatorname{Homeo}(O)}(Q),\id\right)
       \longrightarrow \pi_1(\operatorname{Homeo}(O),\id).
                                                        \tag{C.32.1}
\]
Moreover, a based null-homotopic loop in
$\operatorname{Homeo}^+(\R^2)$ preserving a connected open subset
$O\subseteq\R^2$ restricts to a null-homotopic loop in
$\operatorname{Homeo}(O)$.
Consequently, for a hyperbolic centralizer loop and an open orbit in
C.28 whose closed restricted operator image in C.30 is discrete, its
restriction has a based equivariant null-homotopy on that orbit.
The original restriction kernel is retained; no extension of this
null-homotopy to the orbit frontier is asserted.
\end{claim}
\begin{proof}
First consider the free proper discrete action on $O$. It gives an
actual covering
\[
                       p:O\longrightarrow V=O/Q.          \tag{C.32.2}
\]
Here discreteness refers to the topology in the proper-action
hypothesis. For completeness, a compact neighborhood of a point
meets only finitely many of its translates, by properness. Shrink
an open neighborhood inside it to avoid the nonidentity translates
among this finite set, using freeness. Distinct translates of the
result are disjoint, which gives the evenly covered neighborhoods
for (C.32.2). The quotient $V$ is Hausdorff by properness, is a
surface without boundary by these local charts, and is second
countable because images of a countable open basis of $O$ form an
open basis. The connected covering is regular, with deck group
exactly $Q$.

We record the homeomorphism-group input precisely. Every
second-countable surface admits a PL triangulation. Applied with
the compact subpolyhedron consisting of one point $v$, the compact
case of \cite[Lemma~2.1]{C_Yagasaki2000} and the noncompact case of
\cite[Theorem~1.1]{C_Yagasaki2000} give
\[
 \operatorname{Homeo}(V,v)^0\simeq *
 \quad\text{for }V\not\cong S^2,\ \R P^2,\ \R^2.          \tag{C.32.3}
\]
These are all exceptions relevant to a surface without boundary
with that one-point subpolyhedron. Evaluation
$\operatorname{Homeo}(V)^0\to V$ is a locally trivial bundle:
in a coordinate disc a continuously chosen homeomorphism supported
in that disc moves its center to a nearby point, giving a local
section. Its identity fibre component is exactly the group in
(C.32.3). The homotopy exact sequence therefore gives an injection
\[
 \pi_1(\operatorname{Homeo}(V),\id)
          \longrightarrow\pi_1(V,v)
       \quad(V\not\cong S^2,\R P^2,\R^2).                \tag{C.32.4}
\]
The base $V$ of (C.32.2) cannot be $S^2$ or $\R P^2$:
their fundamental groups are finite, so every connected covering
has finitely many sheets and compact total space, contrary to
noncompactness of $O$. If $V=\R^2$, its connected covering is
one-sheeted, so $Q=1$ and (C.32.1) is the identity homomorphism.

Let $f_t$ be a based loop in the centralizer that is null-homotopic
in $\operatorname{Homeo}(O)$. It descends to a based loop $b_t$
on $V$, defined by the exact equation
\[
                         b_t p=p f_t.                    \tag{C.32.5}
\]
The family is continuous: joint continuity descends through the
open quotient map $[0,1]\times O\to[0,1]\times V$, and the
compact-open exponential law applies since $V$ is locally compact.
For $x\in O$, the evaluation loop $f_t(x)$ is null-homotopic in
$O$. Thus $b_t(p(x))=p(f_t(x))$ is null-homotopic in $V$.
By (C.32.4), $b_t$ is a null-homotopic based loop of homeomorphisms
of $V$. Choose a continuous filling
$B:D^2\to\operatorname{Homeo}(V)^0$ with the prescribed boundary
$b_t$, and a point $z_*\in\partial D^2$ corresponding to the based
endpoint, so that $B_{z_*}=\id$.

The lift is a filling through homeomorphisms and through the exact
centralizer; we verify both assertions. Form
\[
 F:O\times D^2\longrightarrow V,\qquad
                       F(x,z)=B_z(p(x)).                 \tag{C.32.6}
\]
At $(x_*,z_*)$, the induced fundamental-group image is
$p_*\pi_1(O,x_*)$, since $D^2$ is simply connected and
$F(x,z_*)=p(x)$. The covering lifting criterion gives a unique
continuous lift $\widetilde B:O\times D^2\to O$ with
$\widetilde B(x,z_*)=x$.
Apply the identical construction to
$F^-(x,z)=B_z^{-1}(p(x))$, giving a continuous lift
$\widetilde C$ that is the identity at $z_*$.
The maps
$(x,z)\mapsto\widetilde B_z\widetilde C_z(x)$ and
$(x,z)\mapsto x$ lift the same map $(x,z)\mapsto p(x)$
and agree at the basepoint. Uniqueness of lifts makes them equal;
the reverse composition is equal for the same reason.
Hence every $\widetilde B_z$ is a homeomorphism, with inverse
$\widetilde C_z$, and the family is compact-open continuous.

For $q\in Q$, the two maps
$\widetilde B_z q$ and $q\widetilde B_z$ lift the same map
$B_zp$ and agree at $z_*$, so they are equal for every $z$.
On $O\times\partial D^2$, both $\widetilde B_z$ and the specified
$f_t$ lift (C.32.5) and agree at the based endpoint. They are
therefore equal on that entire connected boundary product.
Thus $\widetilde B$ fills the original loop in its centralizer.
This proves (C.32.1).

Now let $O$ be a connected open subset of the plane preserved by
a based ambient-null loop. Its restriction is a continuous loop
in $\operatorname{Homeo}(O)^0$. By
\cite[Theorem~1.1]{C_Yagasaki2000}, with the empty subpolyhedron,
this identity group is homotopically trivial unless $O$ is the
plane or the open annulus; the nonorientable and boundary
exceptions in the source cannot occur for $O\subseteq\R^2$.
In either remaining case its fundamental group is $\Z$.
Choose model coordinates on $O$, a coordinate Jordan circle
$J\subset O$, and two distinct points $a,b$ of $J$. A Schoenflies
homeomorphism $h$ of the ambient plane sends $J$ to the unit circle.
The map
\[
 \Theta:\operatorname{Homeo}(O)\longrightarrow S^1,\qquad
 k\longmapsto
       \frac{h(k(b))-h(k(a))}
            {|h(k(b))-h(k(a))|}
                                                        \tag{C.32.7}
\]
has induced homomorphism $\pi_1\operatorname{Homeo}(O)\to\Z$
that is an isomorphism. Indeed, the full rotation loop in the
chosen plane or annulus model preserves $J$ and has restriction
winding one. The explicit computation in C.29 shows that its
image under (C.32.7) also has winding one. A homomorphism
$\Z\to\Z$ whose image contains $1$ is an isomorphism.
Formula (C.32.7) extends as a continuous map on all ambient
plane homeomorphisms by the same two evaluations. The
ambient-null loop therefore has zero image under it, so its
restricted fundamental-group class is zero. This proves the
second assertion in all cases.

Finally take an open orbit from C.28 and let $Q$ be the closed
restricted operator image supplied by C.30. When this group is
discrete, C.30 gives a free proper action with precisely that
discrete topology, so (C.32.2) applies. The restricted canonical
plane loop commutes with $Q$ and is null-homotopic in
$\operatorname{Homeo}(O)$ by the second assertion. Its
null-homotopy supplied by (C.32.1) commutes with $Q$, hence with
the image of the operator stabilizer. The kernel in (C.30.4)
acts identically on $O$ throughout. This proves the stated
equivariant conclusion with the actual restriction homomorphism.
\end{proof}

\paragraph{Transport and frontier control.}
One may choose the obtained contraction on one representative of
an operator orbit of open $G$-orbits and transport it to the others
by conjugation. Equivariance under the setwise stabilizer makes
this independent of the transporting operator. This establishes
well-defined orbitwise families. The covering lift above supplies
no uniform estimate at orbit frontiers. C.33 subsequently treats
all positive-dimensional proper closed operator images
intrinsically, and C.34 proves joint continuity on the full
nonopen locus. C.35 records the obstruction to extending its
particular median prescription, while C.36 constructs compatible
alternatives on finite unions of closures with two-dimensional
proper operator image. The remaining frontier and accumulation
requirements are stated in C.12, which remains open.
